\section{Regular observability and joint area--jet coordinates}
\label{sec:v14-observability}
We distinguish the information in a normalized germ from that in a
finite set of physical probabilities. The distinction also specifies
which finite-dimensional families the two-flight observation can
recover. The rank criterion below applies to a supplied family; the
subsequent construction verifies that criterion for a physical family
with an unknown area.

\subsection{Exact normalization and finite observations}
\label{sec:v14-normalization}
In Section~\ref{sec:v13-two-flight}, the geometric inverse takes the
\emph{normalized} germs $G_0,G_1$ as data and supplies
$g,\kappa_0,\kappa_1$ separately. To form those germs from raw
probabilities using the definition of $G_b$, $A$ must either be
supplied or recovered. This is automatic only for exact germ data:
with $p_b(d)=\Pr(E_{2,b}(d))$,
\begin{equation}\label{eq:v14-exact-area}
 \ell:=\lim_{d\downarrow0}\frac{p_b(d)}{d^2}
       =\frac{1}{2A\sinh(2\gamma)},\qquad
 A=\frac{1}{2\ell\sinh(2\gamma)},\qquad
 G_b(d)=\frac{p_b(d)}{\ell d^2}.
\end{equation}
The leading coefficient is the same in the two orientations. These
identities follow from $G_b(0)=1$. Extracting $\ell$ from an exact
germ is not a finite-data statistical procedure. In
Theorem~\ref{thm:v13-two-flight-observation}, $A$ is fixed and supplied
as part of the physical family. Theorem~\ref{thm:v14-area-windows}
below removes that supplied area for a specified enlarged family using
only positive windows. Neither statement removes the separate labelled
curvatures from the contact-jet inverse.

\subsection{The regular rank criterion}
Let $\theta$ range over an open subset of $\R^p$. Assume that channel
labels and a common physical measurement schedule are supplied,
independent of the unknown parameter. For instance, at fixed supplied
$g$ the schedule is $2g+d$. Let $p_b(d;\theta)$ denote the actual
two-flight probability, $b=0,1$, $0<d<d_*$. Suppose these functions
are continuously differentiable in $\theta$. Whenever we speak of
analytic germs, assume that they and their first parameter derivatives
are analytic in $d$ on a common neighborhood of zero. Define
\begin{equation}\label{eq:v14-tangent-observation}
 \mathcal L_{\theta_*}v=
 \bigl(d\mapsto D_\theta p_b(d;\theta_*)v\bigr)_{b=0,1}.
\end{equation}

\begin{theorem}[Finite positive windows and regular rank]
\label{thm:v14-rank}
For any fixed open interval $J\subset(0,d_*)$, the following are
equivalent:
\begin{enumerate}
\item no nonzero $v\in\R^p$ annihilates both tangent functions on $J$;
\item there are $p$ labelled positive windows $(b_i,d_i)$, $d_i\in J$,
whose probability map has nonsingular derivative at $\theta_*$;
\item such a $p$-window map is bi-Lipschitz on a sufficiently small
compact parameter ball about $\theta_*$ with nonempty interior.
\end{enumerate}
If the tangent germs are analytic, their linear independence is
equivalent to these conditions on every nonempty open $J$ in their
connected common collar. Thus every regularly observable supplied
finite-dimensional family admits a finite positive-window design;
finite dimension alone is not the hypothesis. The constants and the
chosen windows are local to that family.
\end{theorem}
\begin{proof}
For each $(b,d)$ let $\ell_{b,d}$ be the covector
$D_\theta p_b(d;\theta_*)$. If their span were a proper subspace of
$(\R^p)^*$, its annihilator would contain a nonzero vector, contrary
to the first condition. Hence $p$ of these covectors form a basis,
which proves the second condition. Conversely an invertible evaluation
matrix has zero kernel.

Write $P$ for the selected $p$-window map and $J_*=DP(\theta_*)$.
Continuity permits a convex neighborhood on which
$\|J_*^{-1}DP-I\|<1/2$. Integration along the line segment between
any two parameters gives
\[
 \|J_*^{-1}(P(\theta)-P(\theta'))-(\theta-\theta')\|
                           \le\tfrac12\|\theta-\theta'\|.
\]
Together with a bounded derivative, this proves both Lipschitz bounds.
Restrict to a small closed ball. Conversely, differentiating a
bi-Lipschitz lower bound at its interior point $\theta_*$ proves
injectivity of the derivative, hence nonsingularity in dimension $p$.
Finally, the analytic identity theorem says that a tangent linear
combination vanishing on $J$ vanishes on the whole connected collar
and as a germ. This proves the last assertion.
\end{proof}

\begin{corollary}[Local coordinates for the observable quotient]
\label{cor:v14-observable-quotient}
Suppose the span of $D_\theta p_b(d;\theta)$, over the supplied
windows in a fixed collar, has constant rank $r$ near $\theta_*$.
Locally there are coordinates $(x,z)\in\R^r\times\R^{p-r}$ such
that all selected probabilities depend only on $x$, and $x$ can be
taken to be $r$ positive-window probabilities. Equality of these $r$
means is locally equivalent to equality of the entire selected laws.
This is an exact local factorization, not statistical sufficiency of
finite noisy sample means.
\end{corollary}
\begin{proof}
Choose $r$ independent evaluation covectors at $\theta_*$ and shrink
the neighborhood so that they stay independent. The corresponding
map is a submersion. By the constant-rank assumption its differential
has kernel equal to the common kernel of all selected tangent
observations. A local submersion chart writes this map as $x$. In a sufficiently small
product neighborhood, each scalar probability has zero $z$ derivative,
and is therefore constant on each connected $z$ slice. Thus it is a
function of $x$. The converse follows because the $r$ coordinates
are themselves selected probabilities. The case $r=0$ says that all
selected probabilities are locally constant.
\end{proof}

The hypothesis concerns regular local inversion. A map with a singular
derivative may still be injective with a weaker modulus, as in the
scalar map $\theta\mapsto\theta^3$. The theorem does not exclude
such singular experiments.

The local geometric information set is also essential. Keep the two
selected facing obstacles and the free area fixed, and translate a
third obstacle inside a remote region separated from all the selected
collision and residual-time charts. Its area is unchanged. For a
small translation, the selected stationary actions, twists, and
first-hit domains are unchanged. The exact physical flux integral
therefore gives the same selected probabilities throughout their
common collar, at every flight number. Generic fixed selected
obstacles remove ambient isometries, so this is a nonconstant family
of tables. This example, identified by the referee, is a zero-rank
direction for the selected observation, not a counterexample to
Theorem~\ref{thm:v13-two-flight-observation}. The target of that theorem
is its physical contact-coordinate family, not an arbitrary entire
billiard table.

\subsection{A physical family with a free area coordinate}
\begin{proposition}[Independent area and contact jets]
\label{prop:v14-free-area}
For each fixed $M\ge2$ and the fixed curvature radii in
Theorem~\ref{thm:v12-realization}, its support construction extends to
a real-analytic family of dimension $2M-1$ with local coordinates
\[
                (A,q_2,\ldots,q_M).
\]
The selected gap and the two labelled curvatures stay fixed; the
contacts remain individually even. Each sufficiently small fixed-area
slice retains independent contact jets and the selected leading
hierarchy of Theorem~\ref{thm:v12-realization}.
\end{proposition}
\begin{proof}
In \eqref{eq:v12-support}, keep the coefficient $z$ of
$\sin^{2M+2}\theta$ free instead of solving the area constraint.
To avoid confusion with the hyperbolicity parameter in
Section~\ref{sec:v13-two-flight}, denote this free coefficient by $\zeta$.
It changes neither support values nor the two curvature radii, and
it changes no graph jet through order $2M$, because it vanishes to
order $2M+2$ at both selected normals. The Jacobian of
$s\mapsto(q_2,\ldots,q_M)$ retains the nonzero triangular diagonal
\eqref{eq:v12-support-jacobian}.

For $\phi(\theta)=\sin^{2M+2}\theta$, differentiation of the exact
area formula and integration by parts yield
\[
 \partial_\zeta A
  =-\partial_\zeta\mathcal S(h)
  =-\int_0^{2\pi}(h+h'')\phi\,\dd\theta<0.
\]
Strict positivity of $h+h''$ and nonnegativity of the nonzero function
$\phi$ give the strict inequality. In coordinates $(s,\zeta)$ the
Jacobian of $(q_2,\ldots,q_M,A)$ is consequently block triangular
with invertible diagonal blocks. The analytic inverse-function theorem
provides the claimed coordinates. Positivity of curvature, clearance,
the physical selected channel and evenness persist under the same
small restriction used in the original support construction. Fixing
$A$ recovers its area constraint; in particular the slice
$A=12-\pi$ is the original family.
\end{proof}

\subsection{Joint recovery from a minimal regular set of windows}
\begin{theorem}[Two-flight area--jet observation]
\label{thm:v14-area-windows}
Fix $M\ge2$, put $n=M-1$, and supply the labelled geometry
$g,\kappa_0,\kappa_1$ and the analytic physical family of
Proposition~\ref{prop:v14-free-area}. The area $A$ is unknown.
There are a fixed $h>0$ and a compact coordinate ball $\Theta$ with
nonempty interior on which the $2n+1=2M-1$ actual means
\begin{equation}\label{eq:v14-area-windows}
 P(\theta)=
 \bigl(p_0(ih;\theta)_{1\le i\le n+1},
       p_1(ih;\theta)_{1\le i\le n}\bigr)
\end{equation}
are bi-Lipschitz coordinates for $\theta=(A,q_2,\ldots,q_M)$.
All physical observation times are $2g+ih$. No probability derivative,
zero-offset value, or supplied area enters this design. The number
$2M-1$ is minimal for a differentiable bi-Lipschitz coordinate map
made of scalar window means on this full-dimensional family.

For sufficiently small $\varepsilon>0$ and $0<\eta<1/4$, a Borel
estimator using independent binary preparations satisfies
\begin{equation}\label{eq:v14-area-confidence}
 \sup_{\theta\in\Theta}
 \Pr_\theta\{\|\widehat\theta-\theta\|>\varepsilon\}\le\eta,
 \qquad N_{\rm tot}\le C\varepsilon^{-2}\log(C/\eta).
\end{equation}
For all sufficiently large $N$, the minimax squared risk, allowing
adaptive and randomized selection among these fixed windows and at
most $N$ fresh preparations, lies between $c/N$ and $C/N$.
All preparations, including failures, are counted. Constants may
depend on the fixed family and $M$; no uniform growing-order or
full-profile minimax assertion is made.
\end{theorem}
\begin{proof}
Set $\alpha=(2A\sinh(2\gamma))^{-1}$. The finite-jet theorem and
the common analytic integral give
\begin{equation}\label{eq:v14-scaled-expansion}
 \frac{p_b(d;\theta)}{d^2}
 =\alpha\left(1+\sum_{k=1}^n\xi_{b,k}d^k
                            +d^{n+1}R_b(d;\theta)\right),
\end{equation}
with uniformly bounded first parameter derivatives of the remainders
on a compact restriction. The normalized coefficients through order
$n$ depend only on the fixed leading contact geometry and the displayed
graph jets. Higher jets may vary with $A$ in this family, but their
contributions belong to the remainder.

Use the analytic coordinates
\[
 \beta=(\alpha,\alpha\xi_{0,1},\ldots,\alpha\xi_{0,n},
                  \alpha\xi_{1,1},\ldots,\alpha\xi_{1,n}).
\]
They are legitimate coordinates: the first entry determines $A$,
division by its positive value gives both $\xi$ vectors, and
Theorem~\ref{thm:v13-two-flight} recovers the graph jets.
The leading nodal Jacobian for the scaled probabilities in
\eqref{eq:v14-scaled-expansion} is $V_h$, with rows
\[
 (1,ih,\ldots,(ih)^n,0,\ldots,0)\quad(1\le i\le n+1),
\]
and
\[
 (1,0,\ldots,0,ih,\ldots,(ih)^n)\quad(1\le i\le n).
\]
It factors as
$V_h=V_1\diag(1,h,\ldots,h^n,h,\ldots,h^n)$.
To see that $V_1$ is invertible, let a vector belong to its kernel.
Its first $n+1$ equations say that a polynomial of degree at most
$n$ vanishes at $1,\ldots,n+1$, so its constant and all first-block
coefficients vanish. The remaining polynomial has zero constant and
vanishes at $1,\ldots,n$; its additional root at zero forces all
second-block coefficients to vanish. Consequently
$\|V_h^{-1}\|\le C h^{-n}$ for $0<h\le1$.

The remainder in the nodal Jacobian is $O(h^{n+1})$. Multiplication
by $V_h^{-1}$ makes it $O(h)$. Fix $h$ so small that every node lies
in the common collar and this error has norm less than $1/2$.
Returning from scaled to actual probabilities multiplies rows by
$(ih)^2$, which are positive \emph{known} constants. There is no row
normalization involving the unknown area. The resulting physical
Jacobian is nonsingular. Shrink to a convex coordinate ball where
its product with the inverse central Jacobian stays within $1/2$ of
the identity. The segment argument in Theorem~\ref{thm:v14-rank}
gives the two Lipschitz bounds. Use the compact closure as $\Theta$.
All these finitely many probabilities then belong to
$[p_*,1-p_*]$ for some $p_*>0$, after decreasing the collar if
necessary. A map with fewer than $2n+1$ scalar means cannot have an
injective derivative on a $(2n+1)$-dimensional interior. Differentiating
a bi-Lipschitz lower bound proves the claimed regular minimality.

For clarity, the statistical argument uses these unnormalized physical
means. Write $K=2n+1$, take $L$ independent preparations per window,
and let $Y$ be their sample-mean vector. The bounded-summand exponential
bound gives
\[
 \Pr_\theta\{\|Y-P(\theta)\|_\infty>t\}
                                      \le2K\exp(-2Lt^2).
\]
Choose a finite ordered parameter net whose image has Euclidean mesh
at most $t$. Minimize the discrepancy from $Y$, breaking ties by the
first index. This is a Borel estimator, since it makes finitely many
comparisons of continuous functions of $Y$. Its image error is at most
$2\|Y-P(\theta)\|+t$, and the inverse Lipschitz bound controls its
parameter error. Taking $t=c\varepsilon$ and
$L=\lceil C\varepsilon^{-2}\log(2K/\eta)\rceil$ proves
\eqref{eq:v14-area-confidence}, with deterministic charge $KL$.
For squared risk, take $L=\lfloor N/K\rfloor$ and net mesh
$N^{-1/2}$. The same deterministic inequality and
$\mathbb E\|Y-P(\theta)\|^2\le C/L$ give the upper bound $C/N$.

For the lower bound choose an interior point $\theta_*$ and physical
alternatives $\theta_\pm=\theta_*\pm aN^{-1/2}v$, where $v$ is a
unit vector. The smooth mean map and the positive probability margins
give, uniformly over the allowed windows,
\[
 D(\operatorname{Ber}(p_+)\Vert\operatorname{Ber}(p_-))
 \le\frac{(p_+-p_-)^2}{p_-(1-p_-)}\le\frac{Ca^2}{N}.
\]
Given the past and a procedure's independent random seed, its rule for
choosing the next window is the same under the two parameters. The
conditional entropy chain rule therefore bounds the entropy of the
entire adaptive record by $Ca^2$. An early stop can be padded with
parameter-independent null observations. Choose $a>0$ so small that
the total variation distance is at most $1/2$. The test choosing the
nearer of the two parameters from an estimator has summed error at
least $1/2$. On a test error its squared loss is at least $a^2/N$,
so the maximum risk is at least $a^2/(4N)$. Both alternatives belong
to the actual family of Proposition~\ref{prop:v14-free-area}.
\end{proof}

The same proof applies to any supplied analytic family for which
$(A,q_2,\ldots,q_M)$ is a full local coordinate at fixed labelled
$g,\kappa_0,\kappa_1$, with the stated common integral and remainder
bounds. Proposition~\ref{prop:v14-free-area} verifies those hypotheses
rather than inferring them from dimension alone. The exact germ
normalization \eqref{eq:v14-exact-area}, the joint finite-window
experiment, and the smooth-class pilot in
Section~\ref{sec:v12-regularized-observation} are three different operations.
The last removes $g,A,\gamma$ for the profile experiment under its
supplied class certificates; it is not being used as an uncharged
oracle in the contact-jet theorem.
